% =============================================================================
% compliance.tex — CUMCM Compliance Layer（v4.5 模板 v2 §3.1）
% 只放官方硬约束与最终检查接口；不决定字体审美/标题字号/颜色（那些在 style_tokens.tex）。
% 本文件以注释形式给出"官方法规事实"，供写作与门禁对齐；模板参数不在这里改。
% =============================================================================

% ---------- 官方硬约束（事实层，来自竞赛格式规范/试行规定） ----------
% CUMCM-1  纸张与编译：A4、xelatex（中文字体经字体栈自行回退，禁止本机绝对字体路径）。
% CUMCM-2  摘要页：标题 + 摘要 + 关键词同页（第 1 页，页码 1 页脚居中）；无目录页。
% CUMCM-3  摘要原则上 ≤ 1 页；"600-900 字"是风格建议，不是官方数值（本模板不作硬规则）。
% CUMCM-4  附录必须含参赛建模用到的全部完整可运行源程序（非核心摘录）。
% CUMCM-5  支撑材料作为单个 ZIP 提交；论文附录 A 的清单须与 ZIP 内容一一对应
%          （由 SUBMISSION_MANIFEST 自动生成，禁止硬编码 5 行）。
% CUMCM-6  AI 工具使用声明位于参考文献之前（二选一官方定句；使用版须经参赛队确认后定稿，
%          未确认前只允许"（内容待参赛队确认）"占位，禁止自动断言未使用/只用于某环节）。
% CUMCM-7  参考文献引用须真实可核验（verify_refs 门逐条 OpenAlex/Crossref 核验）。

% ---------- 模板开关（与 SECTION_PLAN.json 对应；见 template_config.tex） ----------
% 本模板的章节显示由 \showdatasection 等开关控制；问题数由 \questioncount 控制。
% 默认关闭的可选内容：第二章框架/第六章稳健性与结论；模型评价保留为第六章父层。
